\section{Correspondence with the formal proof}\label{app:lean}

The formal proof is part of the Lean library of \cite{OpenAI26} (Lean~4.34.1 with Mathlib). The
statement checked by Comparator is
\begin{center}
\small
\texttt{theorem OAI.DirichletCharacter.}\\
\texttt{\ \ LFunction\_ne\_zero\_of\_twenty\_three\_twenty\_fourths\_lt\_re}\\
\texttt{\ \ \{q :\ }$\mathbb N$\texttt{\} [NeZero q] (}$\chi$\texttt{ :\ DirichletCharacter }$\mathbb C$\texttt{ q)}\\
\texttt{\ \ \{s :\ }$\mathbb C$\texttt{\}}
\texttt{ (hs :\ (23 / 24 :\ }$\mathbb R$\texttt{) < s.re) (hpole :\ }$\neg$\texttt{(}$\chi$\texttt{ = 1 }$\wedge$\texttt{ s = 1)) :}\\
\texttt{\ \ DirichletCharacter.LFunction }$\chi$\texttt{ s }$\ne$\texttt{ 0}\\[3pt]
(Comparator configuration \texttt{ComparatorChallenges/\allowbreak DirichletTwentyThreeTwentyFourths.json}),
\end{center}
together with its corollary for $\zeta$, which has no $s\ne1$ hypothesis:
\texttt{OAI.riemannZeta\_ne\_zero\_of\_\allowbreak twenty\_\allowbreak three\_\allowbreak twenty\_\allowbreak fourths\_\allowbreak lt\_\allowbreak re}
(configuration \texttt{Quasi\allowbreak Riemann\allowbreak Twenty\allowbreak Three\allowbreak Twenty\allowbreak Fourths.json}).
Their only axioms are \texttt{propext}, \texttt{Classical.choice} and \texttt{Quot.sound}. The
modules \texttt{OAI.NumberTheory.\allowbreak DirichletL.\allowbreak QRH.*} state the steps of this paper as Lean theorems
in the notation used here: the scale is $D$, and $M\ge bD$ is a cutoff beyond the support of $W$
that makes every sum finite. (Inside the underlying library the roles are reversed: there $Z$ is
the scale and $D\in\mathbb N$ the cutoff.) The main theorem is proved by composing these steps
exactly as in Section~\ref{sec:reduction}, and Proposition~\ref{thm:ms} is proved as
Proposition~\ref{prop:poisson-reduction} applied to Proposition~\ref{prop:canonical}. Names below
are in the namespace \texttt{OAI.QRH}, in the module named in parentheses.

\medskip
{\small\raggedright
\begin{tabular}{@{}p{0.34\textwidth}p{0.62\textwidth}@{}}
\toprule
This paper & Lean\\
\midrule
$A_u(D)$, $A^\varphi(D)$, $S$ & \texttt{A }$\chi$\texttt{ W D M u}, \texttt{weightedSum }$\chi$\texttt{ W D M}, \texttt{S q} (\texttt{Setup})\\
Theorem~\ref{thm:main} & the theorem above (module \texttt{OAI.NumberTheory.\allowbreak DirichletL.\allowbreak QRH})\\
Proposition~\ref{thm:ms} & \texttt{meanSquare} (\texttt{MeanSquare})\\
Proposition~\ref{prop:extraction} & \texttt{saving\_of\_meanSquare} (\texttt{Extraction})\\
Proposition~\ref{prop:mellin} & \texttt{LFunction\_ne\_zero\_of\_saving} (\texttt{ZeroFree})\\
Lemma~\ref{lem:arithmetic} & \texttt{gauss\_jacobi}, \texttt{gauss\_two\_cube}, \texttt{mobius\_absorption} (\texttt{Milestones})\\
Lemma~\ref{lem:poisson} & \texttt{poisson\_with\_exclusions}\\
Proposition~\ref{prop:poisson-reduction} & \texttt{poissonReduction} (\texttt{MeanSquare})\\
Proposition~\ref{prop:canonical} & \texttt{dualBound} (\texttt{MeanSquare}); \texttt{descent\_bound}, \texttt{descent\_by\_rank}\\
Proposition~\ref{prop:R} & \texttt{completed\_mean\_square}\\
Lemma~\ref{lem:cube-reduction} & \texttt{cube\_inversion}, \texttt{short\_cubes}, \texttt{descent\_step}\\
Proposition~\ref{prop:transfer} & \texttt{transfer}; Lemma~\ref{lem:arithmetic-poisson}: \texttt{second\_poisson}\\
Proposition~\ref{lem:reflection} & \texttt{theta\_reflection}; \texttt{thetaModels} (\texttt{MeanSquare})\\
Lemma~\ref{lem:theta-bounds} & \texttt{theta\_cusp\_bounds}\\
Lemma~\ref{lem:quadratic} & \texttt{quadratic\_large\_sieve}\\
Patterson's coefficients \eqref{eq:intro-theta-coefficients}, automorphy of $\theta$ & \texttt{theta\_coefficients}, \texttt{theta\_automorphy}\\
$E_\varphi(s)$, $P_S(s)$, \eqref{eq:weighted-euler} & \texttt{OAI.SmoothMobiusCorrection.\allowbreak eulerCorrection}, \texttt{OAI.CompactMellinBridge.\allowbreak finiteEulerFactor}\\
\bottomrule
\end{tabular}\par}

\medskip
Names without a module are aliases in \texttt{Milestones}, whose documentation gives the
library theorem behind each. The formal proof of
Sections~\ref{sec:initialization}--\ref{sec:transfer-proof} follows the same strategy as the
text: Poisson summation in $u$, completed sums, theta reflection, the quadratic large sieve, cube
inversion, and the transfer by two Poisson summations. It differs in four points of
bookkeeping, recorded in Remarks~\ref{rem:formal-dual} and~\ref{rem:formal-descent}:
\begin{itemize}
\item test functions are expanded in height twists $\NK(\mathfrak n)^{\mathrm i\theta}$ once and
for all, rather than controlled by $C^J$-norms;
\item excluded primes are carried in the dual mean square as a mask, rather than removed by
Lemma~\ref{lem:remove-exclusions};
\item rows are weighted by a compactly supported smooth majorant, whose Fourier transform only
decays rapidly, rather than by a weight with compactly supported Fourier transform;
\item the descent uses a fixed cube cutoff and an explicit loss budget.
\end{itemize}
The import closure of \texttt{OAI.NumberTheory.\allowbreak DirichletL.\allowbreak QRH} contains none of the modules of
the argument of \cite{OpenAI26} (its probe, detector and moment inductions), and no
\texttt{sorry}.
